%%%%%%%%%%%%%%%%%%%%%%%%%%%%%%%%%%%%%%%
% Cover Letter - CesiumAstro
% Embedded Software Engineering Internship (Summer 2027)
% Plain XeLaTeX — matches simple resume style (Spectral font)
%%%%%%%%%%%%%%%%%%%%%%%%%%%%%%%%%%%%%%%

\documentclass[11pt]{article}

\usepackage[margin=0.75in, top=0.7in, bottom=0.7in]{geometry}
\usepackage[hidelinks]{hyperref}
\usepackage{parskip}
\usepackage{fontspec}

\setmainfont{Spectral}[
  Path = ../../cv/fonts/,
  UprightFont    = Spectral-Regular.ttf,
  BoldFont       = Spectral-Bold.ttf,
  ItalicFont     = Spectral-Italic.ttf,
  BoldItalicFont = Spectral-BoldItalic.ttf
]

\setlength{\parindent}{0pt}
\setlength{\parskip}{7pt}

% No hyphenation
\hyphenpenalty=10000
\exhyphenpenalty=10000

\pagestyle{empty}
\urlstyle{same}

\begin{document}

% ── HEADER ───────────────────────────────────────────────────────────────────
\begin{center}
  {\LARGE\textbf{Travis Brown}}\\[3pt]
  \href{mailto:browntravis760@gmail.com}{browntravis760@gmail.com} \textbullet\
  518-796-8703 \textbullet\
  \href{https://www.linkedin.com/in/tbrown5932}{linkedin.com/in/tbrown5932}
\end{center}

\vspace{2pt}
\noindent\rule{\linewidth}{0.5pt}
\vspace{10pt}

% ── DATE ─────────────────────────────────────────────────────────────────────
\begin{flushright}
  \today
\end{flushright}

\vspace{4pt}

% ── BODY ─────────────────────────────────────────────────────────────────────
Dear CesiumAstro Hiring Team,

I am applying for the Summer 2027 Embedded Software Engineering Internship.
I'm a Computer Science student at RIT in the combined bachelor's and
master's program, with an expected graduation date of December 2027, and
most of what I do outside of class is embedded firmware. At GE Vernova, I
built a full I/O subsystem from scratch in C on an
STM32F7 for industrial radios, using the STM32 Hardware Abstraction Layer, and
implemented the Modbus register map that exposes diagnostic data over the
network. On the RIT Electric Vehicle Team (EVT), I write firmware in both
bare-metal C++ and ThreadX RTOS for a custom electric motorcycle we race. My most recent project is creating our real-time vehicle control unit (VCU) with a 300\,Hz control loop across dual CAN buses. In my free time I'm building an autonomous
fixed-wing flight controller on an ESP32S3 running FreeRTOS, with PID
stabilization and GPS waypoint following, on a KiCad PCB I designed myself.

This is the kind of embedded work I want to be doing. CesiumAstro builds
hardware that has to work in space, and I want to help build the software
that runs on it, whether that's drivers, board support, or the
applications on top. I also like that your engineering teams sit close to
manufacturing and see hardware through from concept to delivery. That's
already how I like to work, from designing my own PCBs to writing the
firmware that runs on them.

On EVT, the VCU firmware does more than just run a state machine. It also handles system-wide
fault detection, directs data logging, and parses data from other onboard
subsystems across multiple threads. The flight controller integrates the
Madgwick AHRS library to fuse IMU and magnetometer data into the attitude
estimates that feed the PID loop. Getting any of that right depends on understanding the
hardware underneath it, so I have found reading schematics and datasheets directly when
bringing up new boards to be crucial.

For protocols and interfaces, I have production firmware experience with
the Modbus register map, plus work across CAN/CANopen, SPI, I2C, UART,
ADC, and PWM. On EVT, I also added STM32F446 support to our team's
hardware-agnostic peripheral library, writing the drivers directly from the
microcontroller's datasheets and reference manuals. I debug at the
hardware level regularly, too, using a Saleae logic
analyzer, oscilloscope, and J-Link Segger to trace firmware issues back to
root cause. Writing code that's safe, readable, and easy to build on
matters just as much as getting it working in the first place. At
Lockheed Martin, I worked inside a defense engineering
process built around peer review and defect tracking, and the code review
tool I wrote there was deployed as a mandatory GitLab CI gate on every
merge request.

I would welcome the opportunity to talk more about where I could contribute.
Thank you for your consideration.

\vspace{16pt}
Kind regards,

\vspace{10pt}
Travis Brown

\end{document}
